\section{Poisson 求和公式}
\label{sec:poisson-summation}
% ============================================================================

上一节我们把 Eisenstein 级数写成了
\[
  G_k(\tau) = 2\zeta(k) + 2\sum_{c=1}^{\infty} \sum_{d\in\Z} (c\tau+d)^{-k}.
\]
这已经很像一个答案了，但还不是真正适合计算和比较的答案。
因为这是一种\textbf{按格点 $(c,d)$ 求和}的写法，而我们真正想要的是 $q=e^{2\pi i\tau}$ 的幂级数展开，也就是
\[
  f(\tau)=\sum_{n\ge 0} a_n q^n.
\]
问题就在这里：一个看起来像"二维格点总和"的对象，怎么会突然变成"一维频率展开"？这两种语言表面上几乎毫不相干。

Poisson 求和公式正是连接这两种语言的桥梁。
它告诉我们：\textbf{对整数点求和，可以改写成对 Fourier 频率求和。}
对固定的 $c$，内层和
$\sum_{d\in\Z}(c\tau+d)^{-k}$
本来是沿着整数平移在累加；一旦做 Fourier 变换，它的第 $m$ 个频率会冒出
$\sim m^{k-1}e^{2\pi i mc\tau}$
这样的项，而这正是 $q$-展开里最希望看到的形状。
所以 Poisson 求和不是一个技术小把戏，而是把"格点求和"翻译成"Fourier 系数"的总翻译器。

学会这一步之后，我们立刻获得三样东西。
第一，我们终于能把 $G_k$ 写成显式的 $q$-展开，从而读出它的系数。
第二，我们会看见为什么 Eisenstein 级数的系数自然变成除数和这类算术函数。
第三，你也会顺手碰到现代解析数论里最经典的思想：在原空间里难看的和，到了频域里往往变得整齐。
历史上 Riemann 在 1859 年证明 $\zeta$ 函数的函数方程，本质上用的就是同一套办法——对格上的求和做 Poisson 变换。

Poisson 求和公式是推导 $E_k$ 的 $q$-展开的核心工具，
也是联系格求和与 Fourier 分析的桥梁。
它的思想很简单：对一个函数"按整数周期化"，
有两种方式——直接加格点，或者在频域加频率——结果相同。

\begin{definition}[Fourier 变换]
\label{def:fourier-transform}
\index{Fourier 变换 (Fourier transform)}
设 $f\colon \R \to \C$ 是 Schwartz 函数（光滑且各阶导数迅速衰减）。
其 \textbf{Fourier 变换}定义为
\[
  \hat{f}(\xi) = \int_{-\infty}^{+\infty} f(x)\, e^{-2\pi i \xi x}\, dx.
\]
\end{definition}

\begin{example}[高斯函数的 Fourier 变换]
\label{ex:gaussian-fourier}
\index{高斯函数 (Gaussian function)}
取 $f(x) = e^{-\pi x^2}$（高斯函数，满足 $\hat{f} = f$，
这是 Schwartz 函数中最特别的一个）。直接计算：
\[
  \hat{f}(\xi)
  = \int_{-\infty}^{+\infty} e^{-\pi x^2} e^{-2\pi i \xi x}\, dx.
\]
配方：$-\pi x^2 - 2\pi i \xi x = -\pi(x + i\xi)^2 - \pi\xi^2$，故
\[
  \hat{f}(\xi) = e^{-\pi \xi^2} \int_{-\infty}^{+\infty} e^{-\pi(x+i\xi)^2}\, dx.
\]
通过围道积分（围绕一个矩形区域），积分值等于
$\int_{-\infty}^{+\infty} e^{-\pi x^2}\, dx = 1$（高斯积分）。
所以 $\hat{f}(\xi) = e^{-\pi\xi^2} = f(\xi)$：
\textbf{高斯函数是自己的 Fourier 变换！}

稍作推广：对 $a > 0$，取 $f_a(x) = e^{-\pi a x^2}$，则
$\hat{f}_a(\xi) = a^{-1/2} e^{-\pi \xi^2/a}$。
\end{example}

\begin{theorem}[Poisson 求和公式]
\label{thm:poisson-summation}
\index{Poisson 求和公式 (Poisson summation formula)}
设 $f\colon \R \to \C$ 是 Schwartz 函数（或更一般地，$f$ 和 $\hat{f}$ 均绝对可积且连续）。则
\[
  \sum_{n \in \Z} f(n) = \sum_{n \in \Z} \hat{f}(n).
\]
\end{theorem}

% figures/ch04-poisson-summation.tex
% Poisson 求和公式示意：整数格点求和 ↔ Fourier 频率求和
\begin{figure}[ht]
\centering
\begin{tikzpicture}[>=Stealth, font=\small, scale=1.0]

% === 左侧：f 在整数点的求和 ===
\begin{scope}[xshift=-4cm]
  \draw[->] (-3.2, 0) -- (3.2, 0) node[right] {$x$};
  \draw[->] (0, -0.2) -- (0, 2.5) node[above] {$f(x)$};

  % 连续函数曲线（高斯型）
  \draw[thick, blue!70!black, domain=-3:3, samples=80]
    plot (\x, {2.0*exp(-0.7*\x*\x)});

  % 整数点上的线段（柱高度硬编码）
  % n=-3: h=2e^{-6.3}≈0.04, n=-2: h=2e^{-2.8}≈0.12, n=-1: h=2e^{-0.7}≈0.99
  % n=0: h=2, n=1: h≈0.99, n=2: h≈0.12, n=3: h≈0.04
  \foreach \n/\h in {-3/0.04, -2/0.12, -1/0.99, 0/2.0, 1/0.99, 2/0.12, 3/0.04}{
    \draw[red!80!black, line width=2pt] (\n, 0) -- (\n, \h);
    \fill[red!80!black] (\n, \h) circle (2.5pt);
  }

  \node[below, font=\footnotesize] at (0, -0.5)
    {$\displaystyle\sum_{n \in \Z} f(n)$};
  \node[above right, blue!70!black, font=\footnotesize] at (1.5, 1.6)
    {$f(x)$};
\end{scope}

% === 中间等号 ===
\node[font=\Large] at (0, 1) {$=$};
\node[font=\footnotesize, text=gray!70!black, align=center] at (0, 0.2)
  {Poisson\\求和};

% === 右侧：hat{f} 在整数点的求和 ===
\begin{scope}[xshift=4cm]
  \draw[->] (-3.2, 0) -- (3.2, 0) node[right] {$\xi$};
  \draw[->] (0, -0.2) -- (0, 2.5) node[above] {$\hat{f}(\xi)$};

  % Fourier变换后的高斯（宽度不同）
  \draw[thick, green!50!black, domain=-3:3, samples=80]
    plot (\x, {1.7*exp(-1.4*\x*\x)});

  % 整数点: n=0: 1.7, n=±1: 1.7e^{-1.4}≈0.42, n=±2: 1.7e^{-5.6}≈0.007
  \foreach \n/\h in {-2/0.01, -1/0.42, 0/1.7, 1/0.42, 2/0.01}{
    \draw[orange!80!black, line width=2pt] (\n, 0) -- (\n, \h);
    \fill[orange!80!black] (\n, \h) circle (2.5pt);
  }

  \node[below, font=\footnotesize] at (0, -0.5)
    {$\displaystyle\sum_{n \in \Z} \hat{f}(n)$};
  \node[above right, green!50!black, font=\footnotesize] at (1.2, 1.2)
    {$\hat{f}(\xi)$};
\end{scope}

% === 底部公式 ===
\node[font=\footnotesize, align=center] at (0, -1.4)
  {$\hat{f}(\xi) = \int_{-\infty}^{+\infty} f(x)\,e^{-2\pi i\xi x}\,dx$};

\end{tikzpicture}
\caption{Poisson 求和公式的直觉。对光滑衰减函数 $f$，
在整数格点上的求和 $\sum_{n \in \Z} f(n)$（左，红色竖线）
等于其 Fourier 变换 $\hat{f}$ 在整数点上的求和 $\sum_{n \in \Z} \hat{f}(n)$（右，橙色竖线）。
关键步骤：把 $\sum_n f(n+x)$ 展开为 Fourier 级数，再令 $x=0$。}
\label{fig:poisson-summation}
\end{figure}


\begin{proof}
定义周期函数
\[
  F(x) = \sum_{n \in \Z} f(x + n).
\]
（由 $f$ 的 Schwartz 性质，对每个 $x$ 这个级数绝对收敛，
且 $F$ 是周期为 $1$ 的光滑函数。）

$F$ 有 Fourier 级数展开：
\[
  F(x) = \sum_{m \in \Z} c_m e^{2\pi i mx},
  \qquad c_m = \int_0^1 F(x) e^{-2\pi imx}\, dx.
\]
计算 $c_m$：
\begin{align}
  c_m
  &= \int_0^1 \left(\sum_{n \in \Z} f(x+n)\right) e^{-2\pi imx}\, dx \\
  &= \sum_{n \in \Z} \int_0^1 f(x+n) e^{-2\pi imx}\, dx \\
  &= \sum_{n \in \Z} \int_n^{n+1} f(t) e^{-2\pi im(t-n)}\, dt
    \qquad (t = x + n,\; e^{2\pi imn} = 1) \\
  &= \int_{-\infty}^{+\infty} f(t) e^{-2\pi imt}\, dt = \hat{f}(m).
\end{align}
（交换求和与积分由绝对收敛保证；第三步用了 $e^{2\pi imn} = 1$。）

故 $F(x) = \sum_{m \in \Z} \hat{f}(m) e^{2\pi imx}$。
令 $x = 0$：
\[
  F(0) = \sum_{n \in \Z} f(n) = \sum_{m \in \Z} \hat{f}(m). \qedhere
\]
\end{proof}

\textbf{推广：$t > 0$ 的参数化版本。}
设 $g_t(x) = f(tx)$，则 $\hat{g}_t(\xi) = t^{-1}\hat{f}(\xi/t)$。
应用 Poisson 求和：
\[
  \sum_{n \in \Z} f(nt) = \frac{1}{t} \sum_{n \in \Z} \hat{f}(n/t).
\]
这在推导 $\theta$ 函数的变换律时特别有用。

\begin{example}[高斯求和与 $\theta$ 函数]
\label{ex:theta-poisson}
\index{theta 函数 (theta function)}\index{θ@$\theta(t)$}
取 $f(x) = e^{-\pi x^2}$，$\hat{f}(\xi) = e^{-\pi\xi^2}$。
定义 $\theta(t) = \sum_{n \in \Z} e^{-\pi n^2 t}$（$t > 0$）。
用参数化 Poisson 求和（$f \to f(x)$，$t \to \sqrt{t}$）：
\[
  \theta(t) = \sum_{n \in \Z} e^{-\pi n^2 t}
  = \frac{1}{\sqrt{t}} \sum_{n \in \Z} e^{-\pi n^2/t}
  = \frac{1}{\sqrt{t}} \theta(1/t).
\]
这就是 $\theta$ 函数的\textbf{函数方程}，
也是 Riemann $\zeta$ 函数满足对称函数方程的根源：
\[
  \theta(t) = \frac{1}{\sqrt{t}} \theta(1/t)
  \quad \Longrightarrow \quad
  \xi(s) = \xi(1-s),
  \quad \xi(s) = \frac{1}{2}s(s-1)\pi^{-s/2}\Gamma(s/2)\zeta(s).
\]
\end{example}

\begin{example}[用 Poisson 求和计算 $\sum_{d \in \Z} (c\tau + d)^{-k}$]
\label{ex:poisson-for-eisenstein}
这是推导 $G_k$ 的 $q$-展开的直接预备。
固定 $c \in \Z_{>0}$，$\tau \in \HH$，计算
\[
  S_c(\tau) = \sum_{d \in \Z} \frac{1}{(c\tau + d)^k}.
\]
取 $f(x) = (c\tau + x)^{-k}$，$\hat{f}(\xi) = \int_{-\infty}^{+\infty} (c\tau + x)^{-k} e^{-2\pi i \xi x}\, dx$。

令 $z = c\tau + x$（$\im(c\tau) = c\, \im\tau > 0$），$dx = dz$：
\[
  \hat{f}(\xi)
  = e^{-2\pi i \xi(c\tau - c\tau)} \int_{\im(c\tau) + \R} z^{-k} e^{-2\pi i\xi z}\, dz
  = e^{2\pi i \xi c\tau} \int_{\R + i c\, \im\tau} z^{-k} e^{-2\pi i \xi z}\, dz.
\]

\textbf{分两种情形：}

\textbf{情形 1（$\xi \leq 0$）}：当 $\xi \leq 0$ 时，
$e^{-2\pi i\xi z}$ 在上半平面（$\im z > 0$）没有奇点，
围道可以推到上方无穷远而积分趋于零。
故 $\hat{f}(\xi) = 0$（$\xi \leq 0$）。

\textbf{情形 2（$\xi > 0$）}：$e^{-2\pi i\xi z}$ 在下半平面衰减，
围道推到下方，通过 $z = 0$ 处的极点。由留数定理：
\[
  \int z^{-k} e^{-2\pi i\xi z}\, dz = -2\pi i \cdot \Res_{z=0}\left(z^{-k} e^{-2\pi i\xi z}\right).
\]
在 $z = 0$ 附近展开 $e^{-2\pi i\xi z} = \sum_{j=0}^\infty \frac{(-2\pi i\xi)^j}{j!} z^j$，
则 $z^{-k} e^{-2\pi i\xi z}$ 在 $z^{-1}$ 处的系数（留数）为
$\frac{(-2\pi i\xi)^{k-1}}{(k-1)!}$，故
\[
  \hat{f}(\xi) = e^{2\pi i\xi c\tau} \cdot (-2\pi i) \cdot \frac{(-2\pi i\xi)^{k-1}}{(k-1)!}
  = \frac{(-2\pi i)^k \xi^{k-1}}{(k-1)!} e^{2\pi i\xi c\tau}.
\]

代入 Poisson 求和（注意 $\hat{f}(\xi) = 0$ 对 $\xi \leq 0$，对整数 $\xi = m \geq 1$）：
\[
  S_c(\tau) = \sum_{d \in \Z} f(d) = \sum_{m=1}^{\infty} \hat{f}(m)
  = \frac{(-2\pi i)^k}{(k-1)!} \sum_{m=1}^{\infty} m^{k-1} e^{2\pi i m c\tau}.
\]
\end{example}


% ============================================================================
